\documentclass[11pt]{article}
\usepackage{amsmath,amssymb,amsthm,geometry,hyperref}
\geometry{margin=1in}
\setlength{\emergencystretch}{3em}
\newtheorem{theorem}{Theorem}
\newtheorem{lemma}{Lemma}
\newcommand{\D}{\mathcal D}
\hypersetup{pdftitle={AIM 1048: Band nonoverlap},pdfauthor={AI-assisted research; submitted by Siavash Sadeghi}}
\begin{document}
\title{Band nonoverlap from division by a strip Blaschke factor}
\author{AI-assisted research draft for review}
\date{5 November 2026}
\maketitle
\noindent\textbf{Review status.} This is a proposed resolution of AIM problem
1048, prepared with AI assistance. No human review or formal verification
has taken place. Submitted by Siavash Sadeghi. See the accompanying review
note for scope and provenance.


\begin{theorem}[Graph-bottom water waves]
Let $h>0$, and let $b$ be real, $2\pi$-periodic and $C^1$, with
$\min_x(h-b(x))>0$. For the exact harmonic-extension
Dirichlet--Neumann operator with a flat top and homogeneous Neumann
bottom condition, its ordered Bloch eigenvalues satisfy
\[
 \Lambda_n(\theta_a)\le\Lambda_{n+1}(\theta_b)
 \quad\text{for every }n\ge0\text{ and every }\theta_a,\theta_b\in\mathbb R.
 \tag{1}
\]
Consequently $\max_\theta\Lambda_n(\theta)\le
\min_\theta\Lambda_{n+1}(\theta)$ for every $n$. Eigenvalues are counted
with multiplicity, and touching bands are allowed.
\end{theorem}

The proof first establishes a weighted strip theorem. Its crucial
geometric condition is a positive weight on the middle of a doubled
strip. An arbitrary positive surface weight does not suffice for this
argument.

\section*{1. Weighted strip and its analytic representation}
Fix $P=2\pi$ and $d>0$. Let $p$ be a continuous real periodic function
with $\min p>0$, and let
\[
 \rho_k=\frac{p_k}{\cosh(dk)},\qquad k\in\mathbb Z. \tag{2}
\]
Then $\rho$ is continuous and positive. Indeed it is the value on the
middle line of the harmonic function on $|\operatorname{Im}z|<d$ with
boundary value $p$ on both edges, so $\min p\le\rho\le\max p$ by the
maximum principle. Consider the harmonic extension in the single
strip $-d<t<0$, with zero Neumann data at $t=-d$, and use top mass
$\int_0^P\rho(s)|u(s)|^2ds$. All quasiperiodic functions below have
unitary character $e^{iP\theta}$.
The strip form has the explicit Fourier value
\[
 q_\theta[u]=P\sum_m (m+\theta)\tanh(d(m+\theta))|u_m|^2.
\]
Consequently $q_\theta+\int\rho|u|^2$ is equivalent to the
$H^{1/2}$ norm squared. The form is closed, and the compact embedding
$H^{1/2}\hookrightarrow L^2(\rho ds)$ gives compact resolvent.
Its ordered min--max formula can also be seen directly from an
orthonormal eigenbasis: an $(n+1)$-dimensional subspace contains a
nonzero vector orthogonal to the first $n$ basis vectors, while the
span of the first $n+1$ attains the opposite inequality.

For a top trace
$u(s)=\sum_m u_m e^{i(m+\theta)s}\in H^{1/2}$ define its bottom trace
and holomorphic continuation by
\[
 v(s)=\sum_m\frac{u_m}{\cosh(d(m+\theta))}e^{i(m+\theta)s},
 \qquad
 V(z)=\sum_m\frac{u_m}{\cosh(d(m+\theta))}e^{i(m+\theta)z}.
 \tag{3}
\]
Thus $V$ is holomorphic on $\D_d=\{|\operatorname{Im}z|<d\}$, has
the same character, belongs to $H^1$ on one doubled cell, and its two
edge traces belong to $H^{1/2}$. The original harmonic extension, with
$r=t+d\in[0,d]$, is
\[
 U(s,-d+r)=\frac{V(s+ir)+V(s-ir)}2. \tag{4}
\]
Conversely every holomorphic $H^1$ function on a doubled cell with
unitary character gives an admissible top trace by (4). These claims
follow directly from Fourier series: the multipliers on either edge
in (3) are bounded, and the analytic Dirichlet energy weights a
coefficient $v_m$ by
$(m+\theta)\sinh(2d(m+\theta))$, with its nonnegative limiting
interpretation at zero.

Write $V_\pm(s)=V(s\pm id)$ and $V_0(s)=V(s)$. The strip energy and
weighted top mass are respectively
\[
 Q(V)=\frac12\int_{(0,P)\times(-d,d)}|V'(z)|^2\,ds\,dt, \tag{5}
\]
\[
 M(V)=\frac14\int_0^P\rho(s)
       \bigl(|V_+(s)|^2+|V_-(s)|^2\bigr)\,ds
       +\frac12\int_0^P p(s)|V_0(s)|^2\,ds. \tag{6}
\]
For (5), differentiate (4) and use
$|a+b|^2+|a-b|^2=2|a|^2+2|b|^2$.
For (6), expand the top square in (4). The real part of its cross
term is
\[
 \operatorname{Re}\int_0^P\rho V_+\overline{V_-}\,ds
 =\int_0^P p|V_0|^2\,ds. \tag{7}
\]
For finite Fourier sums this follows because symmetrizing the
coefficient $e^{d((n+\theta)-(m+\theta))}\rho_{n-m}$ replaces it by
$\cosh(d(n-m))\rho_{n-m}=p_{n-m}$. General functions follow by
Fourier truncation: both edge traces converge in $L^2$, and the middle
trace also converges in $L^2$. Notice that (7) is independent of the
Bloch phase.
The real-part operation in (7) is essential: the unsymmetrized
cross integral need not be real for complex data or a non-even
weight. No evenness or reality of the Bloch function is assumed.

\section*{2. An inner factor with any prescribed nontrivial character}
Fix $s_0\in\mathbb R$ and $y\in(-d,d)$, and put
\[
 c=\frac{\pi}{4d},\qquad \tau=\tan(cy),\qquad
 b_y(z)=\frac{\tanh(cz)-i\tau}{1+i\tau\tanh(cz)}.
\]
The map $\tanh(cz)$ maps $\D_d$ conformally to the unit disk;
$b_y$ is its disk Blaschke factor for the point $iy$. Therefore
$|b_y|<1$ inside, $|b_y|=1$ on the two edges, and its only zero in
the strip is $iy$, a simple zero. Its limits as $\operatorname{Re}z$
tends to plus or minus infinity are
\[
 b_+=e^{-2icy},\qquad b_-=-e^{2icy}.
\]
Define
\[
 B_y(z)=b_y(z-s_0)\prod_{j=1}^\infty
 \frac{b_y(z-s_0-jP)}{b_-}
 \frac{b_y(z-s_0+jP)}{b_+}. \tag{8}
\]
Each normalized factor differs from one by an exponentially decaying
quantity as $j\to\infty$, uniformly on a closed period rectangle and
on a small neighborhood of its two horizontal edges. The product and
its derivatives thus converge normally there. It defines a
holomorphic function on the whole strip, smooth through the edges,
whose zeros are exactly $s_0+jP+iy$, all simple. Every factor has
modulus at most one in the closed strip, so
\[
 |B_y|\le1\text{ in }\D_d,
 \qquad |B_y|=1\text{ on its two edges}. \tag{9}
\]
More explicitly, fix $0<\eta<d-|y|$. The nearest poles of a single
factor are at heights $2d-y$ and $-2d-y$, outside
$|\operatorname{Im}z|\le d+\eta$. On a bounded horizontal interval
the normalized tails are $1+O(e^{-2cjP})$, uniformly in that enlarged
strip. Their deviations from one are summable. The infinite tail
therefore has a nonzero limit away from the listed zeros, and normal
convergence on a slightly larger compact set also gives convergence
of derivatives. Thus the assertion about the exact zero set does
not follow merely from boundedness of the finite products.
Reindexing a finite symmetric product and taking its limit gives
\[
 \frac{B_y(z+P)}{B_y(z)}=\frac{b_+}{b_-}
 =\exp\!\left(i\pi\left(1-\frac yd\right)\right).
 \tag{10}
\]
The identity extends through the zeros. Hence its character parameter
is
\[
 \chi(y)=\frac{1-y/d}{2}\in(0,1). \tag{11}
\]
In particular every nontrivial unitary character occurs. Finally
$\log|B_y|$ vanishes on the edges and is negative immediately inside;
its outward normal derivative there is nonnegative. This follows
already by taking the one-sided derivative of a smooth function
that is nonpositive on an inward normal segment and zero at its
endpoint. Write
\[
 \beta_\pm(s)=\partial_\nu\log|B_y(s\pm id)|\ge0. \tag{12}
\]
The product is character-automorphic, not generally periodic.
No uniform edge regularity as $y\to\pm d$ is asserted or needed.

\section*{3. Division decreases energy and increases mass}
For the main argument it suffices to take a finite Bloch Fourier
series $u$, so that $V$ in (3) is entire. Suppose it vanishes at
$a=s_0+iy$. By its character it vanishes at all $a+jP$. Consequently
\[
 W=V/B_y \tag{13}
\]
is holomorphic throughout the strip and has character parameter
$\theta-\chi(y)$ modulo integers. The apparent poles are removable,
and $B_y$ is analytic and nonzero through the two horizontal edges.
Thus $W$ is smooth through the closed doubled cell. It represents
an admissible top trace by (4); no assertion about division of a
rough trace is needed. Reconstruction is injective: on each real
Bloch frequency it multiplies the middle Fourier coefficient by
$\cosh(d(m+\theta-\chi(y)))>0$, after an integer relabeling if needed.

The equal edge moduli and (6), together with $p>0$ and $|B_y|\le1$,
give
\[
 M(V)\le M(W). \tag{14}
\]
For the energy, the exact identity is
\[
 Q(V)-Q(W)
 =\frac14\sum_{\pm}\int_0^P\beta_\pm(s)|W_\pm(s)|^2\,ds\ge0.
 \tag{15}
\]
Indeed, for a holomorphic smooth function $F$,
$\Delta|F|^2=4|F'|^2$, and hence (5) is one eighth of the boundary
integral of $\partial_\nu|F|^2$. The lateral contributions cancel
because $|F|^2$ is periodic. On the horizontal edges $|B_y|=1$ and
\[
 \partial_\nu|B_yW|^2-\partial_\nu|W|^2
 =2\beta_\pm|W|^2.
\]
This proves (15) for precisely the smooth functions used below.

Combining (14)--(15) yields
\[
 \frac{Q(W)}{M(W)}\le\frac{Q(V)}{M(V)}
 \quad(V\ne0). \tag{16}
\]
The division map is complex linear and injective on its domain.

\section*{4. All-phase ordered interlacing on the weighted strip}
Let $\lambda_n(\theta)$ denote the ordered weighted strip
eigenvalues. Fix arbitrary $\theta_a,\theta_b$. If their characters
coincide, $\lambda_n(\theta_a)\le\lambda_{n+1}(\theta_b)$ is the
ordinary ordering. Otherwise choose $y$ so that
\[
 \chi(y)=\theta_b-\theta_a\pmod1,
 \qquad 0<\chi(y)<1. \tag{17}
\]
Fix $\varepsilon>0$. Approximate a basis of the first $n+2$
eigenfunctions in phase $\theta_b$ by finite Bloch Fourier series in
$H^{1/2}$. Since the form and mass Gram matrices converge, with the
mass matrix remaining positive definite, this gives an
$(n+2)$-dimensional finite Fourier subspace $L$ such that
\[
 \max_{0\ne u\in L}\frac{q_{\theta_b}[u]}{\int\rho|u|^2}
 \le\lambda_{n+1}(\theta_b)+\varepsilon. \tag{17a}
\]
This remains true when eigenvalues have multiplicity: any orthonormal
eigenbasis for their span can be approximated in this way.
Evaluation $V\mapsto V(s_0+iy)$ is one complex linear functional on
$L$. Its kernel intersection has dimension at least $n+1$.
Apply division (13) and reconstruct top traces by (4). This gives an
injective linear image of dimension at least $n+1$ in phase
$\theta_a$, and (16) holds for every nonzero vector. The min--max
principle consequently gives
\[
 \lambda_n(\theta_a)\le\lambda_{n+1}(\theta_b)+\varepsilon.
\]
More explicitly, the maximum quotient on $L$ is at least the maximum
on its kernel intersection, which is at least the maximum on the
divided image, which is at least $\lambda_n(\theta_a)$. Letting
$\varepsilon\downarrow0$ proves
\[
 \lambda_n(\theta_a)\le\lambda_{n+1}(\theta_b). \tag{18}
\]
The approximation is only in one finite-dimensional spectral
subspace for fixed $n$ and phase. No uniform-in-index approximation
or division of a general form-domain function is being assumed.

\section*{5. Exact geometric reduction for analytic graph bottoms}
Suppose first that $b$ is real analytic and periodic, with positive
depth $H(x)=h-b(x)$. The physical channel modulo horizontal period
is conformally an annulus. For example $\zeta=e^{-iz}$ identifies
it with the doubly connected planar domain bounded by the unit
circle and the analytic simple curve
$\zeta=e^{-H(x)}e^{-ix}$; its two boundary components are disjoint.
Annular uniformization, with its analytic boundary extension, gives
after lifting logarithms a conformal map
\[
 Z:\{-d<\operatorname{Im}w<0\}\longrightarrow\Omega_b,
 \qquad Z(w+P)=Z(w)+P, \tag{19}
\]
for some $d>0$, taking top to top and bottom to bottom. The map and
its inverse extend analytically locally through either boundary,
and their derivatives are nonzero there. The horizontal period is
still exactly $P$; no depth-to-period ratio is discarded.

For precise standard references, annular uniformization is the
connectivity-two case of L.~V.~Ahlfors, \emph{Complex Analysis},
third edition, McGraw--Hill, 1979, Chapter VI, Section 5.1,
Theorem 10, pp.~255--256. The analytic boundary extension is the
local reflection argument in Chapter VI, Sections 1.3--1.4,
Theorems 3--4, pp.~233--235; locally straightening an analytic
boundary arc and a circular annulus boundary gives the same
reflection situation. The required hypotheses here are analytic
embedded boundary circles separated by a positive distance.

On the top $Z(s)$ is real and strictly increasing, so
\[
 \rho(s)=Z'(s)>0.
\]
On the bottom write $Z(s-id)=X(s)+i(-h+b(X(s)))$. This parametrizes
the graph once with the positive horizontal orientation, so
\[
 p(s)=\operatorname{Re}Z'(s-id)=X'(s)>0. \tag{20}
\]
For completeness, $X'$ cannot vanish: the derivative of a regular
parametrization of a $C^1$ graph is proportional to
$(1,b'(X))$, and the conformal derivative is nonzero. Orientation
and degree one fix its positive sign.
The periodic holomorphic function $Z'$ is real on the top. Its
Fourier series therefore has conjugate-symmetric coefficients,
and taking its real part at depth $d$ gives
\[
 p_k=\cosh(dk)\rho_k, \tag{21}
\]
exactly (2). Both weights are continuous and strictly positive.

Conformal invariance of the two-dimensional Dirichlet integral
transforms the physical harmonic minimization into the flat-strip
minimization. The natural bottom condition remains zero Neumann.
The physical surface measure becomes $dx=\rho(s)ds$, and (19)
preserves the Bloch character exactly. The image of a conformal
period cell may have curved lateral sides; the energy and mass
densities are periodic, so integration over this fundamental cell
agrees with integration over the physical vertical-sided cell.
Thus the physical spectrum
is precisely the weighted strip spectrum used in (18). This proves
(1) for analytic graph bottoms, without amplitude, parity,
fundamental-period or genericity assumptions.

\section*{6. Passage to every admissible $C^1$ graph}
Choose real trigonometric polynomials $b_j\to b$ in $C^1$ and keep
a common positive lower depth. With $H_j=h-b_j$ and $H=h-b$,
pull the physical bulk energies back by $y=H_j(x)z$ and $y=H(x)z$
to the fixed cell $(0,P)\times(-1,0)$. Their coefficient matrices
are
\[
 A_H(x,z)=
 \begin{pmatrix}
 H&-zH'\\
 -zH'&(1+z^2(H')^2)/H
 \end{pmatrix}, \tag{22}
\]
acting on $((\partial_x+i\theta)U,U_z)$. Uniform positivity of
the depths and $C^1$ convergence imply uniform ellipticity and
uniform convergence of these matrices. Consequently, for numbers
$\varepsilon_j\to0$,
\[
 (1-\varepsilon_j)q_{b,\theta}
 \le q_{b_j,\theta}\le
 (1+\varepsilon_j)q_{b,\theta}, \tag{23}
\]
uniformly in phase. This follows first for every bulk competitor
and then by taking the infimum with its fixed top trace. Min--max
therefore gives convergence of every fixed ordered eigenvalue,
including its multiplicities. For fixed $n,\theta_a,\theta_b$,
pass (1) for $b_j$ to the limit using (23). This proves (1) for
$b$. Since the phase pair was arbitrary, it gives the asserted
max--min inequality.

For clarity, phase continuity follows directly from that same bulk
realization. If $u$ is the top trace of $U$ and $d_+=\max H$, vertical
integration gives
\[
 \|U\|_{L^2(\mathcal C_b)}^2
 \le 2d_+\|u\|_{L^2}^2+2d_+^2\|U_y\|_{L^2(\mathcal C_b)}^2.
\]
Writing $a_\theta[U]$ for the bulk energy, expansion of
$|(\partial_x+i\theta+i\delta)U|^2$ and Cauchy--Schwarz therefore
give, for $|\delta|\le1$,
\[
 |a_{\theta+\delta}[U]-a_\theta[U]|
 \le C|\delta|\bigl(a_\theta[U]+\|u\|_{L^2}^2\bigr),
\]
where $C$ depends only on $d_+$. Apply the upper bound to the
minimizing extension at $\theta$, and interchange the two phases
for the reverse bound. The resulting two-sided comparison of
$q_\theta+\|u\|^2$ and $q_{\theta+\delta}+\|u\|^2$, followed by
min--max, proves continuity of each fixed ordered eigenvalue.
Compactness of the phase interval gives its maximum and minimum.

The functional-analytic ingredients are the representation of a
densely defined closed nonnegative form by its self-adjoint
operator and its compact-resolvent min--max principle; the form
domain here is $H^{1/2}$ on the top circle and its embedding in
$L^2$ is compact. Standard form background is T.~Kato,
\emph{Perturbation Theory for Linear Operators}, second edition,
corrected printing, Springer, 1980, Chapter VI, Theorem 2.1.
All phase-changing comparisons in the proof were established
explicitly by a single analytic evaluation constraint and inner
factor division. No finite-rank assertion about a vertical
cross-section, or oscillation theorem for real Bloch functions,
is used.
\end{document}
